%%% BEGIN fig:dy-abelian-model
\begin{figure}[htbp]
\centering
\begin{tikzpicture}[fd,
  fc lab/.style={font=\small},
  fc note/.style={font=\footnotesize},
  fc amp/.style={hadron blob,ellipse,minimum width=0.95cm,minimum height=1.9cm,inner sep=0pt,font=\footnotesize}]
  % =====================================================================
  % (a) n exchanges between the two eikonal spectator lines
  % =====================================================================
  \begin{scope}
    \node[fc lab] at (-3.15,0) {$\displaystyle\sum_{n\ge0}$};
    \draw[wilson=0.1] (-2.5,0.75) -- (1.9,0.75);
    \draw[wilson=0.1] (-2.5,-0.75) -- (1.9,-0.75);
    \foreach \x in {-1.35,-0.55,1.05} {
      \draw[glauber] (\x,-0.75) -- (\x,0.75);
      \node[vertex] at (\x,0.75) {};  \node[vertex] at (\x,-0.75) {}; }
    \node[fc lab] at (0.25,0) {$\cdots$};
    \node[fc note,anchor=south west] at (-2.5,0.86) {$s_A$：沿$n_A$，荷$g_A$};
    \node[fc note,anchor=north west] at (-2.5,-0.86) {$s_B$：沿$n_B$，荷$g_B$};
    \draw[gray!60!black,line width=0.5pt,{Stealth[length=3.5pt,width=3pt]}-{Stealth[length=3.5pt,width=3pt]}] (2.2,0.7) -- (2.2,-0.7);
    \node[fc lab,anchor=west] at (2.25,0) {$\boldsymbol b$};
    \draw[line width=0.5pt] (-1.45,-1.43) -- (-1.45,-1.55) -- (1.15,-1.55) -- (1.15,-1.43) (-0.15,-1.55) -- (-0.15,-1.65);
    \node[fc note] at (-0.15,-1.93) {$n$次交换（含全部交叉排列）};
    \node[fc lab,anchor=west] at (2.85,0.3) {$\displaystyle=\sum_{n\ge0}\frac{\big(\ii\chi(\boldsymbol b)\big)^n}{n!}=\ee^{\ii\chi(\boldsymbol b)}$};
    \node[fc lab,anchor=west] at (2.85,-0.75) {$\displaystyle\chi(\boldsymbol b)=-\frac{g_Ag_B}{2\pi}\,K_0(m|\boldsymbol b|)$};
    \node[fc lab] at (2.3,-2.6) {(a)\ 程函旁观者之间$n$次有质量阿贝尔玻色子交换之和为相位};
  \end{scope}
  % =====================================================================
  % (b) the two cut placements at O(g^2)
  % =====================================================================
  \begin{scope}[yshift=-4.95cm]
    \begin{scope}[xshift=-1.2cm]
      \draw[wilson=0.76] (-1.6,0.5) -- (1.6,0.5);
      \draw[wilson=0.76] (-1.6,-0.5) -- (1.6,-0.5);
      \node[fc amp] at (-1.6,0) {$\mathcal M^{(0)}$};
      \node[fc amp] at (1.6,0) {$\mathcal M^{(0)*}$};
      \draw[glauber] (-0.62,-0.5) -- (-0.62,0.5);
      \node[vertex] at (-0.62,0.5) {};  \node[vertex] at (-0.62,-0.5) {};
      \draw[cut] (0,-1.15) -- (0,1.15);
      \node[fc note] at (-0.62,0.78) {$s_A$};  \node[fc note] at (-0.62,-0.8) {$s_B$};
      \node[fc lab] at (0,-1.55) {$+\ii\chi(\boldsymbol b)\,|\tilde{\mathcal M}^{(0)}(\boldsymbol b)|^2$};
    \end{scope}
    \node[fc lab] at (1.55,0) {$+$};
    \begin{scope}[xshift=4.3cm]
      \draw[wilson=0.3] (-1.6,0.5) -- (1.6,0.5);
      \draw[wilson=0.3] (-1.6,-0.5) -- (1.6,-0.5);
      \node[fc amp] at (-1.6,0) {$\mathcal M^{(0)}$};
      \node[fc amp] at (1.6,0) {$\mathcal M^{(0)*}$};
      \draw[glauber] (0.62,-0.5) -- (0.62,0.5);
      \node[vertex] at (0.62,0.5) {};  \node[vertex] at (0.62,-0.5) {};
      \draw[cut] (0,-1.15) -- (0,1.15);
      \node[fc note] at (0.62,0.78) {$s_A$};  \node[fc note] at (0.62,-0.8) {$s_B$};
      \node[fc lab] at (0,-1.55) {$-\ii\chi(\boldsymbol b)\,|\tilde{\mathcal M}^{(0)}(\boldsymbol b)|^2$};
    \end{scope}
    \node[fc lab,anchor=west] at (6.55,0) {$=0$};
    \node[fc lab] at (2.3,-2.3) {$\order{g^{2n}}$：$\displaystyle\sum_{j=0}^{n}\frac{(\ii\chi)^j}{j!}\,\frac{(-\ii\chi)^{n-j}}{(n-j)!}=0\ (n\ge1)$；全阶：$|\ee^{\ii\chi(\boldsymbol b)}|^2=1$};
    \node[fc lab] at (2.3,-3.2) {(b)\ 割线位置之和（对旁观者相对横动量$\boldsymbol\Delta$积分，即在$\boldsymbol b$空间中）};
  \end{scope}
\end{tikzpicture}
\caption{阿贝尔模型中旁观者之间的 Glauber 交换（\cref{prop:dy-abelian-model}）。两个旁观者$s_A$、$s_B$近似为沿$n_A$、$n_B$（方向相反）的无限长程函线，荷为$g_A$、$g_B$，横向间距为$\boldsymbol b$；暗红色螺线为质量为$m$的阿贝尔矢量玻色子（Feynman 规范）；时间自左向右，双线上的箭头为荷的流向。
(a)~对$n$次交换（含全部交叉排列）求和给出冲量参数空间的相位：$n$次项为$(\ii\chi)^n/n!$，总和为$\ee^{\ii\chi(\boldsymbol b)}$，$\chi(\boldsymbol b)=-\frac{g_Ag_B}{2\pi}K_0(m|\boldsymbol b|)$为实数（$m\to0$时红外发散，但始终是相位）。
(b)~$\order{g^2}$的两个割图：交换位于割线左侧（振幅）给出$+\ii\chi$，位于右侧（共轭振幅）给出$-\ii\chi$；对旁观者相对横动量$\boldsymbol\Delta$单举的截面中（Parseval 恒等式化为$\boldsymbol b$空间）二者逐阶相消，全阶时$|\ee^{\ii\chi}|^2=1$，但对固定的$\boldsymbol\Delta$逐点并不相消（VF-DY-13，VF-GLA-3）。}
\label{fig:dy-abelian-model}
\end{figure}
%%% END fig:dy-abelian-model
